\section{Related work}\label{sec:related}

\paragraph{Resource-aware noninterference and relational cost} Ngo et al.\ and RelCost
are the closest work and are discussed in \cref{sec:intro-classical,sec:classical}.
Ngo et al.\ extend automatic amortised resource
analysis~\cite{hofmann2003static,hoffmann2012multivariate}, and relational cost analysis
has since been extended to functional-imperative programs~\cite{Qu2019ARel}. Relative to
Ngo et al., our contribution is the theorem that indexing by size cannot be sound for
opaque operations whose cost depends on content. Relative to RelCost, whose refinements
can express the content comparison for a fixed pair of runs, it is the resolution over
feasible secret outcomes, the observers, the leakage bound, and the application to LLM
calls.

\paragraph{Timing channels, constant time and cross-copying} Agat's transformation pads
the arms of a secret branch so that they take equal time~\cite{agat2000transforming}.
The program-counter security model of Molnar et al.\ asks that the sequence of executed
instructions not depend on secrets, and their transformation removes secret-dependent
branches~\cite{molnar2006program}. Constant-time programming forbids secret-dependent
control flow and memory access altogether; it has been checked by
self-composition~\cite{BacelarAlmeida2013SelfComp} and by product
programs~\cite{Almeida2016ConstantTime}. Hermes, a
reversible language for lightweight encryption published in this journal, uses secret and
public types to make operations on secret data independent of it and so to prevent
timing attacks~\cite{Mogensen2022Hermes}. In all of these the observable operation is
fixed by the program: an instruction, a memory access, a step of time. For an LLM call,
the provider fixes the cost from the request content, so ``the same operation'' has to
mean the same request. OrchLang accepts secret branches whose arms send identical
requests and bounds what the remaining differences reveal.

\paragraph{Relational verification} Noninterference is a two-run
property~\cite{goguen1982security}, and self-composition reduces it to a property of one
program~\cite{Barthe2004SelfComposition,Terauchi2005Safety}; product programs connect
such reductions to relational program logics~\cite{Barthe2016ProductPrograms}.
Probabilistic relational Hoare logic~\cite{barthe2009formal} and coupling
proofs~\cite{barthe2017coupling} relate the randomness of two runs. \Cref{lem:equal} is a
derivation of this kind in which every call is related by the identity coupling, and
comparing signatures is a self-composition that never builds the product program.

\paragraph{Quantitative information flow} Min-entropy leakage~\cite{smith2009foundations},
its behaviour under cascading~\cite{espinoza2012min}, bounds by the number of possible
observations~\cite{kopf2010vulnerability,doychev2013cacheaudit}, and the partition model
of adaptive attacks~\cite{kopf2007information} are what \cref{thm:leakage} uses.
CacheAudit bounds cache leakage by counting observations~\cite{doychev2013cacheaudit}; we
count signature classes, because the provider's noise makes counting observations
useless.

\paragraph{Information-flow typing and declassification} OrchLang's flow typing is
conventional: the lattice model of Denning~\cite{denning1976lattice}, the type-based
formulation of Volpano et al.~\cite{volpano1996sound}, and the standard two-point
confidentiality and integrity lattices surveyed by Sabelfeld and
Myers~\cite{sabelfeld2003language}. Its declassification is trusted, justified in writing
and all or nothing: in the dimensions of Sabelfeld and Sands~\cite{SabelfeldSands2009} it
constrains only where in the program a release happens, not what is released or by
whom, and it offers no robustness guarantee~\cite{Myers2006Robust}. In this
journal, Manzino and de Latorre embed a DSL in Haskell whose types enforce delimited
release, so that declassification cannot reveal more than
intended~\cite{manzino2026haskell}, and Ben Said and Cristescu verify noninterference of
web-service compositions and generate secure orchestrators from
them~\cite{BenSaid2020Orchestration}. Both control which data reach which party; neither
addresses what the pattern of an orchestrator's requests reveals.

\paragraph{Static cost analysis and certified bounds} Static cost analysis goes back to
Wegbreit~\cite{Wegbreit1975}; COSTA derives cost relations from bytecode and solves them
into closed-form bounds~\cite{Albert2012CostBytecode}, and Ali, Lamo and Pun apply the
same method to Rpl, a resource-sensitive workflow modelling language, in this
journal~\cite{ali2023cost}. In Rpl the cost of a step is declared by the program, so costs
compose; in an LLM workflow a tokenizer decides it, and token counts do not compose
(\cref{sec:bounds}). Proof-carrying code~\cite{Necula1997PCC} and resource bound
certification~\cite{CraryWeirich2000} attach checkable evidence to a program. OrchLang's
certificate is a report rather than such evidence, and a checker for it is future work.

\paragraph{Languages for LLM programs} LMQL treats prompting as programming, with
constraints on generation~\cite{BeurerKellner2023LMQL}, and DSPy compiles declarative
pipelines of model calls and optimises their prompts~\cite{Khattab2024DSPy}. Both are
more expressive than OrchLang, and neither checks information flow or certifies a token
bound. OrchLang gives up expressiveness, including the agents that make up most of our
exclusions, in exchange for decidable analyses.

\paragraph{Security of LLM agents} Indirect prompt injection~\cite{greshake2023not} is what
the integrity part of OrchLang's typing targets. CaMeL separates a planner from untrusted
data and tracks provenance and capabilities in a custom
interpreter~\cite{debenedetti2026defeating}; FIDES tracks confidentiality and integrity
labels in an agent planner~\cite{costa2025securing}; Beurer-Kellner et al.\ catalogue
design patterns for agents, among them the plan-then-execute pattern that our AgentDojo
ports follow~\cite{beurerkellner2025design}. AgentFlow combines a flow-centric policy
language with a runtime monitor, and explicitly leaves timing, token-length and
resource-usage channels out of scope~\cite{ammanaghattashivakumar2026agentflow}; Ghost in
the Agent tracks information flow through the execution traces of
agents~\cite{cai2026ghost}; Agentproof verifies extracted workflow graphs
statically~\cite{xavier2026agentproof}. None of these bounds resource use or reasons
about what the content of requests reveals. AgentDojo~\cite{debenedetti2024agentdojo}
supplies the tasks and recorded attacks of our evaluation.

\paragraph{Side channels of LLM services} The token lengths of streamed responses leak
response text to a network observer~\cite{weiss2024what}; output token counts leak
information about inputs through timing~\cite{zhang2024time}; the sizes and timing of
encrypted streams reveal prompt topics across 28 models~\cite{mcdonald2025whisper}; and
prompt caches shared across users form a timing channel at seven
providers~\cite{gu2025auditing}. These papers establish the channel for single calls. We
study how a workflow's control flow feeds it.

\paragraph{Budgets and tokenization} Token Budgets catalogues 63 budget-overrun incidents
in LLM agents and enforces spending caps at run time, with affine types that prevent a
budget from being copied or spent twice~\cite{khan2026token}; its over-reservation
figures are measured on a different corpus and are not comparable with our slack.
Tokenizers are known to charge languages very unequally~\cite{petrov2023language}; we
measure what this means for sound bounds.
